\documentclass[12pt,a4paper]{report}
\usepackage[]{amsmath,amsthm,amssymb,amscd}
\usepackage[latin1]{inputenc}
%\usepackage{fullpage}

\topmargin 0pt
\advance \topmargin by -\headheight
\advance \topmargin by -\headsep
     
\textheight 246mm
     
\oddsidemargin 0pt
\evensidemargin \oddsidemargin
\marginparwidth 0.5in
     
\textwidth 159mm


%               Theorem Environments
\theoremstyle{definition}
\newtheorem{ex}{}
\newcommand{\pd}[1]{\frac{\partial}{\partial #1}} %\pd{x}
%               Subquestions numbered with letters
\renewcommand{\theenumi}{\textbf{\alph{enumi}}}

%               Maths definitions

\newcommand{\NN}{\ensuremath{\mathbb N}}
\newcommand{\ZZ}{\ensuremath{\mathbb Z}}
\newcommand{\CC}{\ensuremath{\mathbb C}}
\newcommand{\RR}{\ensuremath{\mathbb R}}
\newcommand{\TT}{\ensuremath{\mathbb T}}
\newcommand{\g}{\ensuremath{\frak{g}}}
\newcommand{\h}{\ensuremath{\frak{h}}}
\newcommand{\ft}{\ensuremath{\frak{t}}}
\newcommand{\cB}{\mathcal{B}}
\newcommand{\cN}{\mathcal{N}}
\newcommand{\cL}{\mathcal{L}}
\newcommand{\cD}{\mathcal{D}}
%\newcommand{\pd}[1]{\frac{\partial}{\partial #1}} %\pd{x}
%\newcommand{\pd}[1]{\frac{\partial}{\partial #1}} %\pd{x}

\begin{document}
\noindent
{\sffamily\large Marco Zambon 
\hfill   a discutir:   Jueves, 19/04/2012\\Geometr\'ia III, curso 2011-12}
 

\noindent
\hrulefill{}\\
\begin{center}\bfseries\large\textsf{Hoja 4    \vspace{-2mm}}
\end{center}

\noindent
%\hrulefill{}\\
 
 

\begin{ex}
Considera $$M:=\{(x_1,x_2,x_3)\in \RR^3: x_1^2+x_2^2-x_3=0\}.$$

a) Demuestra que $M$ es una subvariedad de $\RR^3$.

b) Encuentra bases de $T_pM$ y $T_qM$,
donde $p:=(0,0,0)\in M$ y $q:=(1,3,10)\in M$.\\
\textbf{Observaci\'on:} Por a) sabemos que $M$  es una subvariedad de  $\RR^3$, as\'i que  $T_qM\subset T_q\RR^3\cong \RR^3$. Una base de $T_q\RR^3$ es dada por $\{\pd{x_1}|_q,\pd{x_2}|_q,\pd{x_3}|_q\}$ (corresponde a la base  canonica de  $\RR^3$ bajo la identificaci\'on $T_q\RR^3\cong \RR^3$). Por lo tanto 
 cada elemento de $T_qM$ es de la forma $a\pd{x_1}|_q+b\pd{x_2}|_q+c\pd{x_3}|_q$ donde $a,b,c\in \RR$. 
\end{ex}


 
\begin{ex} [\textbf{Teorema de la funci\'on inversa}] Demuestra:
sea $f \colon M \to N$ una aplicaci\'on diferenciable entre variedades diferenciales, y sea $p\in M$, entonces:
$$\text{$f_*(p) \colon T_pM \to T_{f(p)}N$ es un isomorfismo $\Leftrightarrow$ $f$ es un difeomorfismo local en $p$}.$$
\end{ex} 






\begin{ex} 
Demuestra que
$$SL(n;\RR) := \{A \in Mat(n;\RR) : det(A) = 1\}$$
es una subvariedad de $Mat(n;\RR)$. >Cual es su dimensi\'on? 
\end{ex} 


\begin{ex} {\bf (Por escrito)} Para qu\'e
 $ c\in \RR$ es $g^{-1}(c)$ una subvariedad de $\RR^2$, donde $$g : \RR^2 \to \RR; (x, y) \mapsto xy?$$  
\end{ex} 

\begin{ex} Sean $M,N$  variedades diferenciales y $f \colon M \to N$ una aplicaci\'on
diferenciable. Demuestra que $$G\colon M \to M \times N, x\mapsto (x,f(x))$$
es un embebimiento.  
\end{ex} 

\begin{ex} 
Considera el campo vectorial $X=\sum_{i=1}^n x_i \frac{\partial}{\partial x_i}$ en $\RR^n$.\\
Cual es el flujo de $X$?  
\end{ex} 


 \end{document}
 %%%%%%% 
 
 

 \begin{ex}
\begin{enumerate}
\item 
\end{enumerate}
\end{ex}
 
 
 
 